\chapter*{How to Read This Book}
\addcontentsline{toc}{chapter}{How to Read This Book}
\markboth{How to Read This Book}{}

\section*{Three Reading Levels}

Each chapter has three levels, distinguished by color. The colors are simply editorial signposts:

\begin{itemize}[itemsep=4pt]
\item \textbf{Main Story (unmarked text)}: Requires only middle school mathematics and common knowledge. Read it through without skipping steps to understand the central causal links. This is the book's backbone.
\item \textbf{[Conceptual Bridge] Green boxes}: A little proportion, graph reading, simple algebra, and probability turns a qualitative ``this is how it works'' into a quantitative ``this is how much.'' A few extra minutes takes your understanding deeper.
\item \textbf{[Further Challenge] Purple boxes}: Windows onto higher levels: orbitals, free energy, reaction mechanisms, and statistical models. Skip a whole box on your first reading without disrupting the main story, and return whenever curiosity calls.
\end{itemize}

\section*{The Route Through Each Chapter}

Each chapter picks up something nearby---a spoonful of salt, an egg, a leaf---and follows a fixed route deeper. First observe the phenomenon, then enter the particle world, look for its rules, and finally introduce the symbols. At the chapter's end, we return to the opening object and explain it with our new model. If what once seemed obvious suddenly becomes full of meaning, the chapter has succeeded.

Old friends also return across chapters: a slice of bread, an apple, a battery, a bar of soap, a pill, a leaf, a drop of blood, a bacterium. They reappear in chapters on elements, bonds, metabolism, and genes, each time revealing another layer. When an old friend returns, ask yourself: how deeply did I understand it last time?

\section*{Five Questions Throughout the Book}

For any material or living process, we ask five questions in turn:

\begin{enumerate}[itemsep=2pt]
\item What is it made of?
\item How are its components connected and arranged?
\item What rules of energy and probability determine whether it can change?
\item How fast does change occur, and how is it controlled?
\item In living things, how is the process preserved, copied, regulated, and selected?
\end{enumerate}

After reading, we hope you can ask these questions yourself about a new medicine in the news, a drink on a shelf, or a fever in your body. Then the book's method will truly be yours.

\section*{What to Try and What Not to Try}

The Try It Yourself features use only things safely obtainable from a kitchen, pharmacy, or stationery shop. Every experiment includes safety instructions; follow them. Other phenomena---alkali metals reacting with water, radioactive decay, or culturing bacteria---are described or illustrated with video resources only, with an explicit instruction not to attempt them at home. Curiosity deserves encouragement, but a good scientist first learns to protect themselves and others.
